\chapter{Модели и алгоритмы оценки доступности городских сервисов с учетом транспортной связности и уличной морфологии}
\label{ch:models-algorithms}

\section{Формализация пространственной структуры города}
\label{sec:city-structure-formalization}

\subsection{Пространственные единицы анализа и агрегация спроса}
\label{subsec:spatial-units-demand}

\subsection{Сервисные объекты, емкости и зоны обслуживания}
\label{subsec:service-capacity-zones}

\subsection{Улично-транспортный граф и интермодальные связи}
\label{subsec:street-pt-graph}

\section{Модель компонент доступности}
\label{sec:accessibility-components-model}

\subsection{Компонента home-to-stop}
\label{subsec:home-to-stop}

\subsection{Компонента stop-to-service}
\label{subsec:stop-to-service}

\subsection{Компонента travel-time по сети общественного транспорта}
\label{subsec:pt-travel-time}

\subsection{Классификация типов провала доступности}
\label{subsec:access-failure-types}

\section{Метод учета street pattern в анализе доступности}
\label{sec:street-pattern-method}

\subsection{Привязка морфологических классов к пространственным единицам}
\label{subsec:morphology-to-units}

\subsection{Home-side, service-side и path-side street pattern}
\label{subsec:street-pattern-sides}

\subsection{Диагностические показатели связи морфологии и провалов доступности}
\label{subsec:morphology-diagnostics}

\section{Модель обеспеченности и неудовлетворенного спроса}
\label{sec:provision-model}

\subsection{Матрица достижимости и временные пороги доступности}
\label{subsec:accessibility-matrix-thresholds}

\subsection{Оценка покрытого и непокрытого спроса}
\label{subsec:covered-uncovered-demand}

\subsection{Практическая интерпретация показателей demand\_left и demand\_without}
\label{subsec:demand-left-without}

\section{Алгоритмы сценарного улучшения доступности}
\label{sec:accessibility-improvement-algorithms}

\subsection{Расширение существующих сервисов}
\label{subsec:service-expansion}

\subsection{Открытие новых сервисных объектов}
\label{subsec:new-service-opening}

\subsection{Транспортные интервенции и генерация дополнительных связей}
\label{subsec:transport-interventions}

\subsection{Сравнение сервисных и транспортных сценариев}
\label{subsec:scenario-comparison}

\section{Выводы по главе}
\label{sec:chapter2-conclusions}
